\begin{textsample}{POS Dim 1 – df – Score 29.00 – df/df\_ef\_2.txt}  \label{ex:f1_pos_100}
Linguagens é \textbf{uma} das áreas do conhecimento \textbf{que} se estende, principalmente, à produção de sentidos na perspectiva de representar o mundo e socializar pensamentos.
De acordo com os Parâmetros Curriculares Nacionais ( 2001 ), o uso de diferentes linguagens, ao longo da história, tem sua importância e valor diretamente relacionados com demandas sociais e culturais de cada momento ( BRASIL, 2001 ).
Nesse sentido, sendo a escola \textbf{um} espaço cuja função precípua é a de democratizar saberes, é importante considerar \textbf{que} o trabalho com as linguagens no Ensino Fundamental pressupõe a articulação entre Língua Portuguesa, Arte ( Dança, Teatro, Música e Artes Visuais ), Educação Física e Língua Estrangeira.
Essa articulação permite a continuidade das experiências \textbf{vividas} na Educação Infantil, expressas em manifestações artísticas, corporais e linguísticas, transitando -as progressivamente para o Ensino Fundamental sem \textbf{que} os objetivos de aprendizagem e conteúdos de cada \textbf{um} dos componentes curriculares se ocultem, mas \textbf{que} se apresentem como parte de \textbf{um} todo com sentido e coerência em relação à vida dos estudantes.
Para o desenvolvimento das linguagens, pressupõe -se leitura relativa à interação do ser humano em suas relações, ao mundo do trabalho e da tecnologia, à produção artística, às atividades de cultura e prática corporal, à área da saúde, aos movimentos sociais, e ainda incorporam saberes como os \textbf{que} advêm das formas diversas de exercício da cidadania, da experiência docente, do cotidiano e dos diversos interesses dos estudantes, na perspectiva de sua formação integral.
As linguagens permitem ao estudante \textbf{uma} leitura mais ampla do meio em \textbf{que} \textbf{vive}, de sua identidade nesse lugar, de quem é o outro como também das relações interpessoais entre os seres humanos.
Elas possibilitam comunicação, \textbf{que} pode ser “ verbal ( oral ou \textbf{visual-motora}, como Libras, e escrita ), corporal, visual, sonora e, contemporaneamente, digital ” ( BRASIL, 2017, p. 61 ), e \textbf{permeiam} todas as atividades humanas na produção de sentidos \textbf{que} representem o mundo e \textbf{que} socializem pensamentos.
Tais atividades permitem a interação das pessoas, constituindo -se como \textbf{sujeitos} sociais e históricos, dotados de conhecimentos, atitudes e valores culturais, morais e éticos.
Com relação à Língua Portuguesa, cabe ressaltar \textbf{que} seu ensino tem passado por diferentes enfoques ao longo da história : do \textbf{método} sintético ao analítico, da \textbf{simples} reprodução mecânica em \textbf{que} o estudante apenas recebe, sem a oportunidade de construir para si o conhecimento, ao caminho para a emancipação, no qual o estudante tem a oportunidade de pensar, compreender e reconstruir, sendo \textbf{um} \textbf{sujeito} ativo no processo de aprendizagem.
Tendo em vista \textbf{que} a língua é \textbf{um} instrumento de poder, pois, por meio dela, efetiva -se a comunicação, construção de conhecimentos, apropriação dos meios científicos, tecnológicos, participação em processos políticos e expressão cultural, é responsabilidade da escola garantir a todos os estudantes acesso a saberes construídos \textbf{historicamente} pela \textbf{humanidade} em relação à língua.
Nesse sentido, ressalta -se \textbf{que} a finalidade precípua do ensino da Língua Portuguesa é propiciar aos estudantes a competência comunicativa, ou seja, a capacidade de expressar -se adequadamente em qualquer situação, de forma oral e escrita, portanto, ler e escrever proficientemente de modo a “ [ … ] resolver problemas da vida cotidiana, ter acesso aos bens culturais e alcançar participação plena no mundo letrado ” ( BRASIL, 2001, p.41 ).
Nesse contexto, ampliar a competência comunicativa de estudantes, pensando na participação social, pressupõe o ensino da Língua Portuguesa por meio de textos concretizados nos mais variados gêneros e suportes \textbf{que} circulam na sociedade, cumprindo funções específicas de comunicação ( ANTUNES, 2009 ).
A partir desse ensino \textbf{que} contemple o trabalho didático com gêneros textuais, é possível a articulação entre oralidade, leitura/escuta, escrita/produção textual e análise linguística/semiótica, pois saberes provenientes de cada \textbf{uma} dessas práticas de linguagem se relacionam na compreensão e utilização de diferentes gêneros textuais, diversificando e ampliando situações de letramento vivenciadas por estudantes.
Marcuschi ( 2008, p.149 ) confirma essa perspectiva de ensino da Língua Portuguesa ao \textbf{dizer} \textbf{que} “ [ … ] o trato dos gêneros \textbf{diz} respeito ao trato da língua em seu cotidiano nas mais diversas formas ”.
Dell’Isola ( 2007 ) considera \textbf{que} gêneros textuais são vias de acesso ao letramento e propõe \textbf{que} o ensino da língua se dê por meio de textos encontrados na vida diária, ou seja, carregados de sentidos, levando -se em consideração a heterogeneidade de textos existentes em nossa sociedade e a necessidade de tornar os estudantes proficientes leitores e produtores de texto.
Assim, é importante \textbf{que} o professor entenda \textbf{que} gêneros textuais se referem a textos específicos \textbf{que} são encontrados no cotidiano ( poemas, cartas, e-mails, receitas, anúncios, WhatsApp, Twitter, Instagram, vlog, podcast, trailer ), \textbf{enquanto} os tipos textuais \textbf{dizem} respeito a modos textuais ( narração, \textbf{exposição}, injunção/instrução, descrição, argumentação ) \textbf{que} podem aparecer com certa predominância ou articulados entre si na organização interna dos gêneros ( MARCUSCHI, 2008 ).
Nesse contexto, o desenvolvimento da oralidade, leitura/escuta, escrita/produção textual e aprofundamento de análise linguística/semiótica e trato com a literatura se dará por meio do trabalho com gêneros textuais \textbf{que} oportunizem situações em \textbf{que} estudantes tenham contato sistemático, em contextos significativos, com a variedade de gêneros textuais \textbf{que} transitam no meio social.
Já o ensino e a aprendizagem de Língua Estrangeira – LE, na etapa do Ensino Fundamental, têm como propósito o desenvolvimento do \textbf{educando} para a construção do exercício de sua cidadania de forma consciente, participativa e dentro de \textbf{uma} perspectiva de valorização e respeito à diversidade humana.
Além disso, o processo educativo \textbf{que} o currículo de Língua Estrangeira propõe visa à preparação do estudante para as relações no mundo do trabalho e em esferas acadêmicas, tendo como base a visão de \textbf{que} o ensino de outras línguas contribui para o aprimoramento pessoal, a formação ética e o desenvolvimento da autonomia intelectual e do pensamento crítico.
Nesse contexto, e tendo em vista as \textbf{abordagens} \textbf{contemporâneas} de ensino de línguas, as quais apontam para a necessidade da inclusão de dimensões \textbf{historicamente} \textbf{relegadas} à marginalidade do processo em ensinos centrados no caráter metalinguístico, o currículo de LE procura incentivar \textbf{um} processo de ensino/aprendizagem com \textbf{centralidade} no sentido.
Para isso, o trabalho a partir de temas, o \textbf{estímulo} ao protagonismo estudantil e o desenvolvimento da pessoa em \textbf{uma} perspectiva de educação integral são elementos de \textbf{suma} importância para \textbf{que} esse processo se realize de forma consistente e significativa.
Entretanto, o foco nesses aspectos não \textbf{implicam} a \textbf{exclusão} de outras dimensões privilegiadas em \textbf{abordagens} \textbf{que} mantêm a língua como centro.
O redimensionamento dos elementos, implícito em propostas \textbf{contemporâneas}, busca ajustes mais equilibrados, cuidado \textbf{que} tradicionalmente não integra a \textbf{abordagem} majoritariamente instalada nos sistemas de ensino : a gramatical.
Assim, este currículo propõe indicar \textbf{uma} \textbf{abordagem} de caráter mais comunicacional como parte das estratégias de superação das dificuldades \textbf{que} ainda geram resultados indesejados no ensino de LE no Brasil.
A matriz curricular de LE apresenta, então, \textbf{uma} alternativa aberta e plural na oferta linguística.
No entanto, para \textbf{que} se efetive esta proposta na implementação do Currículo do Ensino Fundamental, pressupõem -se outras ações do poder público : formação continuada institucionalmente estabelecida ; investimento nos diferentes níveis formais de apoio à implementação do currículo ; reestruturação de elementos condicionantes da oferta de LE ; e \textbf{estímulo} frequente ao \textbf{compromisso} profissional de professores na experimentação de novas \textbf{abordagens}.
A Arte é \textbf{um} componente curricular, dentro da área Linguagens, capaz de promover diálogos \textbf{que} extrapolam as linguagens oral e escrita, além de contribuir para a formação integral do indivíduo por meio da \textbf{dialética} existente entre a subjetividade e o repertório cultural, seja individual ou social.
No ensino da Arte, o contato do estudante com as diversas linguagens artísticas ( Artes Visuais, Dança, Música e Teatro ) propicia a leitura de mundo e de sua realidade, de forma reflexiva e crítica.
Nesse contexto, esse componente curricular permite a relação do estudante com o contexto social por meio da experiência e do entendimento estético, articulados à compreensão \textbf{histórico-cultural}, a fim de compreender a arte como fenômeno humano.
Pretende -se assim \textbf{que} as diversas manifestações da arte e da cultura formem \textbf{um} indivíduo plural, capaz de conhecer a história construída pela \textbf{humanidade}, o patrimônio do mundo e de se comunicar de forma criativa e sensível a fim de \textbf{que} se fortaleça laços de identidade.
O ensino de Arte foi incluído no currículo escolar da Educação Básica a partir da Lei de Diretrizes e Bases – LDB de 1961 ( Lei 4.024/61 ), fortalecendo -se como componente curricular a partir da LDB de 1996 ( Lei Nº 9.394, de 20 de dezembro de 1996 ), em seu artigo 26, § 2º : “ O ensino da arte constituirá componente curricular obrigatório nos diversos níveis da educação básica, de forma a promover o desenvolvimento cultural dos alunos ” ( BRASIL, 1996, s.p. ).
Em 2016, a Lei Nº 13.278/2016 incluiu as Artes Visuais, a Dança, a Música e o Teatro nos currículos dos diversos segmentos da Educação Básica.
Desta maneira, os sistemas de ensino deverão promover a formação de professores para implantar esses componentes curriculares na Educação Infantil, Ensino Fundamental e Médio, no prazo de cinco anos ( BRASIL, 2016 ).
Para a organização do trabalho pedagógico nos Anos Iniciais, sugere -se \textbf{uma} \textbf{abordagem} integrada das linguagens artísticas, por entender \textbf{que} o professor, \textbf{enquanto} \textbf{um} organizador do espaço social educativo ( VIGOTSKI, 2003 ) tem maior flexibilidade e condições de garantir \textbf{um} trabalho interdisciplinar da arte com as demais linguagens e tecnologias.
Já nos Anos Finais, os conteúdos e objetivos de aprendizagem pautam -se na cronologia histórica aliada à apreensão espiralada das manifestações artísticas próprias de matrizes culturais africanas, orientais e de povos originários, procurando articular -se aos conteúdos dos demais componentes curriculares com vistas ao trabalho interdisciplinar.
Dessa forma, procura -se evitar ou reforçar visões mais particularizadas geograficamente em movimentos artísticos, considerando \textbf{que} seja abordada de maneira integrada, fundamentada e consistente.
Contempla -se ainda os diversos elementos das linguagens artísticas ( Visuais, Dança, Música e Teatro ), contextualizando -os no momento histórico, em \textbf{uma} concepção pedagógica \textbf{que} propicie a formação integral do estudante.
Em relação à Educação Física, esse componente tem como objeto de ensino as manifestações da cultura corporal, \textbf{que} contribui para a formação integral do ser humano, desde seu ingresso na escola, por meio de brinquedo, de jogo simbólico, de movimentos gerais vivenciados mediante atividades orientadas, de iniciação das danças, de ginásticas e de jogos pré-desportivos, entre outras atividades \textbf{que}, ao oportunizar as aprendizagens, favoreçam o desenvolvimento do estudante.
O enfoque dessa \textbf{abordagem} é mais abrangente à medida \textbf{que} valoriza e considera aspectos \textbf{sócio-históricos} de cada atividade trabalhada, como também o contexto em \textbf{que} os estudantes estão inseridos e as aprendizagens motoras individuais, independentemente do nível de habilidades \textbf{que} apresentem.
Assim, é fundamental para formulação de propostas para a Educação Física Escolar localizar “ [ … ] em cada \textbf{uma} dessas manifestações ( jogo, esporte, dança, ginástica e luta ) seus benefícios fisiológicos e psicológicos e suas possibilidades de utilização de instrumentos de comunicação, expressão, lazer e cultura ( … ) ” ( BRASIL,1997, s.p. ).
Entende -se \textbf{que} a Educação Física trata do conhecimento produzido e reproduzido pela sociedade a respeito do corpo e do movimento como \textbf{um} veículo de expressão de sentimentos, como possibilidade de promoção, recuperação, programação e manutenção de \textbf{uma} vida de qualidade.
Assim, a área de Linguagens tem o principal objetivo de possibilitar aos estudantes a participação em práticas de linguagem diversificadas \textbf{que} lhes permitam ampliar conhecimentos e capacidades \textbf{expressivas} em manifestações artísticas, corporais e linguísticas, considerando o processo de constante transformação social.

% matched lemmas: abordagem, centralidade, compromisso, contemporâneo, dialético, dizer, educar, enquanto, estímulo, exclusão, exposição, expressivo, historicamente, histórico-cultural, humanidade, implicar, método, permear, que, relegar, simples, sujeito, sumo, sócio-histórico, um, uma, visual-motor, viver
\end{textsample}
